\section{标量环的变换}

记号约定:$A,B,C,A_{1},A_{2}$是环,$\rho\in\Hom_{\mathsf{Ring}}\left(B,A\right),\rho^{\prime}\in\Hom_{\mathsf{Ring}}\left(C,B\right)$.

\begin{definition}{标量限制}{}
	设$E$是左$A$-模.定义$B$在$E$上的左作用为$\left(\beta,x\right)\mapsto\rho\left(\beta\right)x$,则$E$是左$B$-模.记该左$B$-模为$\rho_{\ast}\left(B\right)$.
\end{definition}

\begin{remark}{}{}
	我们称$\rho_{\ast}\left(B\right)$为``把标量环从$A$限制到$B$得到的模''.
\end{remark}

\begin{proposition}{标量限制保持子模的商模}{}
	设$E$是$A$-模,$M$是$E$的$A$-子模,则$\rho_{\ast}\left(M\right)$是$\rho_{\ast}\left(E\right)$的$B$-子模,并且$\rho_{\ast}\left(E/M\right)=\rho_{\ast}\left(E\right)/\rho_{\ast}\left(M\right)$.
\end{proposition}

\begin{proposition}{标量限制保持线性映射}{}
	设$E,F$是$A$-模,则$u\in\Hom_{A}\left(E,F\right)$是$B$-线性映射.进一步,记$u$在从$A$到$B$的标量限制下得到的线性映射为$\rho_{\ast}\left(f\right)$,则我们有典范$\Z$-模同态$\Hom_{A}\left(E,F\right)\hookrightarrow\Hom_{B}\left(\rho_{\ast}\left(E\right),\rho_{\ast}\left(F\right)\right),u\mapsto\rho_{\ast}\left(u\right)$.
\end{proposition}

\begin{remark}{}{}
	上述典范嵌入不一定是双射.换言之,一个$B$-线性映射$\rho_{\ast}\left(E\right)\to\rho_{\ast}\left(F\right)$不一定是$A$-线性映射.例如,当$A$是除环,$B$是$A$的子除环时,若$B\neq A$,则$B$-向量空间$\rho_{\ast}\left(A_{s}\right)$不是$A$-子空间.
\end{remark}

\begin{proposition}{标量限制与直积,直和都可交换}{}
	设$\left(E_{i}\right)_{i\in I}$是一族$A$-模,则$\rho_{\ast}\left(\prod\limits_{i\in I}E_{i}\right)=\prod\limits_{i\in I}\rho_{\ast}\left(E_{i}\right),\rho_{\ast}\left(\bigboxplus\limits_{i\in I}M_{i}\right)=\bigboxplus\limits_{i\in I}\rho_{\ast}\left(M_{i}\right)$.
\end{proposition}

\begin{proposition}{标量限制反向保持生成集}{}
	设$E$是$A$-模,$S$是$\rho_{\ast}\left(E\right)$的生成集,则$S$也是$E$的生成集.
\end{proposition}

\begin{proposition}{}{}
	设$E,F$是$A$-模.
	\begin{enumerate}
		\item 若$\rho$是满射,则典范映射$\Hom_{A}\left(E,F\right)\hookrightarrow\Hom_{B}\left(\rho_{\ast}\left(E\right),\rho_{\ast}\left(F\right)\right)$是双射,$\rho_{\ast}\left(E\right)$的$B$-子模都是$E$的$A$-子模,$E$的生成集都是$\rho_{\ast}\left(E\right)$的生成集.
		\item 若$\rho$是单射,则$E$的$A$-自由子集都是$\rho_{\ast}\left(E\right)$的$B$-自由子集.
	\end{enumerate}
\end{proposition}

\begin{remark}{}{}
	即使$\rho$是单射,$\rho_{\ast}\left(E\right)$的自由集通常不是$E$的自由子集.例如,对于$\R$的子集$\left\{1,\sqrt{2}\right\}$,这是$\Q$-自由集,也是$\R$-相关集.
\end{remark}

\begin{proposition}{生成集,自由族,基在标量限制下的传递性}{}
	设$E$是$A$-模,$\mathcal{A}\coloneq\left(\alpha_{\lambda}\right)_{\lambda\in L}$是$\rho_{\ast}\left(A_{s}\right)$的元素族,$\mathcal{A}^{\prime}\coloneq\left(a_{\mu}\right)_{\mu\in M}$是$E$的元素族,$\mathcal{A}^{\prime\prime}\coloneq\left(\alpha_{\lambda}a_{\mu}\right)_{\lambda\in L,\mu\in M}$是$\rho_{\ast}\left(E\right)$的元素族.
	\begin{enumerate}
		\item 若$\mathcal{A}$是$\rho_{\ast}\left(A_{s}\right)$的生成系,$\mathcal{A}^{\prime}$是$E$的生成系,则$\mathcal{A}^{\prime\prime}$是$\rho_{\ast}\left(E\right)$的生成系.
		\item 设$\rho$是单射.若$\mathcal{A}$是$\rho_{\ast}\left(A_{s}\right)$的自由族,$\mathcal{A}^{\prime}$是$E$的自由族,则$\mathcal{A}^{\prime\prime}$是$\rho_{\ast}\left(E\right)$的自由族.
		\item 设$\rho$是单射.若$\mathcal{A}$是$\rho_{\ast}\left(A_{s}\right)$的基,$\mathcal{A}^{\prime}$是$E$的基,则$\mathcal{A}^{\prime\prime}$是$\rho_{\ast}\left(E\right)$的基.
	\end{enumerate}
\end{proposition}

\begin{corollary}{有限生成在标量限制下的传递性}{}
	设$\rho_{\ast}\left(A_{s}\right)$是有限生成$S$-模,$E$是有限生成$A$-模,则$\rho_{\ast}\left(E\right)$是有限生成$S$-模.
\end{corollary}

\begin{proposition}{标量限制的函子性}{}
	设$E$是$A$-模,$\rho^{\prime\prime}=\rho\circ\rho^{\prime}$.
	\begin{enumerate}
		\item $\rho_{\ast}^{\prime\prime}\left(M\right)=\rho_{\ast}^{\prime}\left(\rho_{\ast}\left(E\right)\right)$.
		\item 若$\rho\in\Isom_{\mathsf{Ring}}\left(B,A\right),\tilde{\rho}=\rho^{-1}$,则$E=\tilde{\rho}_{\ast}\left(\rho_{\ast}\left(E\right)\right)$.
	\end{enumerate}
\end{proposition}

\begin{remark}{}{}
	设$K$是除环,$E$是$K$-向量空间,$A$是$K$的子环,对$K$的诸有限元素族$\left(\xi_{i}\right)_{i=1}^{n}$,存在$\gamma\in A\setminus\left\{0\right\}$使得$\left(\gamma\xi_{i}\right)_{i=1}^{n}$是$A$的元素族.在此条件下,$E$的每个自由族都是$\varphi_{\ast}\left(E\right)$的自由族,其中$\varphi:A\hookrightarrow K$是典范嵌入.
\end{remark}

\begin{remark}{}{}
	设$E$是$A$-模.记$\mathfrak{b}=\ker\left(\rho\right)$.若$\rho$是满射,则$\mathfrak{b}\subseteq\Ann_{B}\left(\rho_{\ast}\left(E\right)\right)$,而$E$是$\rho_{\ast}\left(E\right)/\mathfrak{b}$-模.
\end{remark}

\begin{definition}{$\rho$-半线性映射}{}
	设$E$是$A$-模,$F$是$B$-模.我们称$\Hom_{B}\left(F,\rho_{\ast}\left(E\right)\right)$的元素为$F$到$E$的(左)$\rho$-半线性映射.
\end{definition}

\begin{proposition}{半线性映射的复合与拟}{}
	设$E$是$A$-模,$F$是$F$-模,$G$是$C$-模.
	\begin{enumerate}
		\item 若$u\in\Hom_{C}\left(G,\rho_{\ast}^{\prime}\left(F\right)\right),v\in\Hom_{B}\left(F,\rho_{\ast}\left(E\right)\right)$,则$v\circ u\in\Hom_{C}\left(G,\left(\rho\circ\rho^{\prime}\right)_{\ast}\left(E\right)\right)$.
		\item 若$\rho\in\Isom_{\mathsf{Ring}}\left(B,A\right),u\in\Hom_{B}\left(F,\rho_{\ast}\left(E\right)\right)$为双射,则$u^{-1}\in\Hom_{A}\left(E,\rho_{\ast}^{-1}\left(F\right)\right)$.
	\end{enumerate}
\end{proposition}

\begin{example}{}{}
	设$E$是$A$-模,$\alpha\in A^{\times}$,则位似$h_{\alpha}:E\to E$是关于$\Ad_{\alpha}:A\to A$的$\rho$-半线性映射.
\end{example}

\begin{remark}{}{}
	设$A=A_{1}\times A_{2},e_{1}=\left(1,0\right)\in A,e_{2}=\left(0,1\right)\in A,E$是$A$-模.一方面,
	\begin{enumerate}
		\item $E=E_{1}\boxplus E_{2}$,其中$E_{1}=e_{1}E,E_{2}=e_{2}E$.
		\item $E_{1}$被$e_{2}$零化,$E_{2}$被$e_{1}$零化,因此$E_{1}$是典范$A_{1}$-模,$E_{2}$是典范$A_{2}$-模.
		\item 对$E$的每个子模$M$都有$M=M_{1}\boxplus M_{2}$,其中$M_{1}=e_{1}M\subseteq E_{1},M_{2}=e_{2}M\subseteq E_{2}$.
	\end{enumerate}
	另一方面,设$F_{1},F_{2}$分别是$A_{1}$-模和$A_{2}$-模,映射$p:A\twoheadrightarrow A_{1}$和$p^{\prime}:A\twoheadrightarrow A_{2}$是典范投影,则
	\begin{align*}
		E=p_{\ast}\left(F_{1}\right)\boxplus p_{\ast}^{\prime}\left(F_{2}\right),\text{其中}F_{1}=e_{1}E,F_{2}=e_{2}E.
	\end{align*}
	因此,$E$是$A$-模.综上,考察$A$上的模$\iff$分别独立地考察$A_{1}$-模和$A_{2}$-模.
\end{remark}











